% =============================================================================
% OPTIONAL REPRODUCIBILITY CHECKLIST - ENGLISH
% Thesis-oriented adaptation of the NeurIPS Paper Checklist format:
% Yes / No / N/A plus evidence or a short justification.
% Source format: https://neurips.cc/public/guides/PaperChecklist
% This is not the official NeurIPS form and does not imply conference compliance.
% Disable it in configuration.tex with:
% \providecommand{\ISISReproducibilityChecklistChoice}{disabled}
% =============================================================================

\chapter{Reproducibility Checklist}\label{app:reproducibility-checklist}

\begin{keypoint}[Optional and applicability-aware]
This checklist is an optional thesis aid. Keep only the items relevant to the
work and mark \textbf{N/A} when an item is not applicable or the requested
information is not available. For each answer, cite a thesis section, appendix,
repository file, dataset record, or provide a brief justification. The format
is adapted from the NeurIPS paper checklist; it is not an official conference
submission form.
\end{keypoint}

\noindent\textbf{Suggested use.}
Complete the checklist after the references and any technical appendix. Review
it with the supervisor before submission, and remove confidential links or
information that should not appear in the public thesis.

\begin{reproducibilitychecklist}

\reproitem{Claims and scope}
{Do the abstract and introduction accurately describe the contribution, scope,
assumptions, and level of generality supported by the results?}
{Point to the abstract, contribution list, and the sections containing the main evidence.}

\reproitem{Limitations}
{Are important limitations, strong assumptions, expected failure modes, and
boundaries of validity discussed?}
{Reference a limitations or threats-to-validity section and explain any omitted limitations.}

\reproitem{Theory, assumptions, and proofs}
{When theoretical results are included, are all assumptions stated and are
complete proofs or proof references provided?}
{Mark N/A for work without theoretical claims; otherwise cite theorem and appendix locations.}

\reproitem{Reproducibility path}
{Is there a practical route for another researcher to reproduce or independently
verify the main contribution?}
{Identify the appendix, repository, model access route, protocol, or other verification mechanism.}

\reproitem{Access to code and data}
{When experiments were performed, are code, data, model artefacts, and execution
instructions available, or is any restriction explained?}
{Give persistent links or repository paths, versions, access conditions, and the reason for any restriction.}

\reproitem{Experimental setting}
{Are datasets, splits, preprocessing, baselines, hyperparameters, selection
procedures, random seeds, stopping rules, and software versions specified?}
{Cite the methodology, appendix table, configuration files, and environment or lock file.}

\reproitem{Uncertainty and statistical evidence}
{Are repeated runs, error bars, confidence intervals, statistical tests, or
other appropriate uncertainty information reported and correctly defined?}
{State what varies across runs, the number of repetitions, and how uncertainty was calculated.}

\reproitem{Compute resources}
{Are the hardware type, memory, run time, storage, and approximate total compute
needed for the reported experiments documented?}
{Reference the environment description and distinguish reported runs from exploratory compute when relevant.}

\reproitem{Research ethics}
{Has the work been checked against applicable research-integrity, safety,
privacy, and professional-ethics requirements?}
{Cite the relevant declaration, review procedure, institutional rule, or explain why the item is not applicable.}

\reproitem{Societal impact}
{Where relevant, are potential harmful uses, fairness concerns, privacy risks,
security risks, and likely deployment consequences discussed?}
{Reference the impact or risk discussion and any mitigation measures.}

\reproitem{Safeguards}
{For artefacts with meaningful misuse or dual-use risk, are proportionate
release safeguards, usage conditions, or access controls described?}
{Mark N/A when no elevated release risk exists; otherwise describe the safeguard and residual risk.}

\reproitem{Licences and terms of use}
{Are third-party datasets, code, models, and other assets properly cited, versioned,
and used in accordance with their licences and terms?}
{List the licence or terms for each important asset and identify any unresolved uncertainty.}

\reproitem{Documentation of new artefacts}
{If the thesis releases new code, data, models, or benchmarks, are they documented
with purpose, contents, provenance, limitations, licence, and maintenance information?}
{Mark N/A when no new artefact is released; otherwise cite the README, data card, or model card.}

\reproitem{Human participants or crowdsourcing}
{If people contributed data or participated in a study, are recruitment,
instructions, consent, compensation, and study procedures documented?}
{Mark N/A when no human participants or crowd workers were involved.}

\reproitem{Ethical or institutional approval}
{When required, were participant risks assessed and approval or exemption
obtained from the appropriate ethics or institutional review process?}
{Provide the non-confidential approval reference, or explain why approval was not required.}

\reproitem{Use of generative AI or LLMs}
{When generative AI or an LLM is an important, original, or non-standard part of
the research method, is its role, version, configuration, and verification described?}
{Mark N/A for ordinary writing or formatting assistance that does not affect the research method or evidence.}

\end{reproducibilitychecklist}

\begin{keypoint}[Final review]
A ``No'' answer is not automatically a defect: it should be accompanied by a
clear explanation and, where possible, a mitigation or future-action note.
Remove placeholder answers and verify that every cited section, file, and link
exists in the submitted version.
\end{keypoint}
